% Hafta 4 — Sinyal işleyicide ne yapılır?
% Derleme: py -3.12 tools/latex/derle.py docs/week-4/assets/h04-12-sinyal.tex
\documentclass[tikz,border=8pt]{standalone}
\usepackage{cen429-cizim}
\begin{document}
\begin{tikzpicture}[node distance=10mm and 24mm]
  \node[baslik] (b) at (0,0) {Sinyal işleyicide ne yapılır?};
  \node[altbaslik] at (b.south west) {sinyal, kesilen kodun ortasında gelebilir};

  \node[kotukutu=58mm, below=16mm of b.south west, anchor=north west] (yanlis) {%
    \kb{YANLIŞ}\\[2mm]
    printf, malloc, free\\
    kilit alan her şey\\[3mm]
    yeniden girişte iç yapı bozulur\\
    $\rightarrow$ kilitlenme / çökme};
  \node[iyikutu=58mm, right=of yanlis] (dogru) {%
    \kb{DOĞRU}\\[2mm]
    yalnız bayrak kur:\\
    \kod{volatile sig\_atomic\_t}\\[3mm]
    asıl işi ana döngü yapar\\
    (\kod{write()} gibi güvenli çağrılar)};

  \node[fit=(yanlis)(dogru), inner sep=0pt] (grup) {};
  \node[dip, text width=150mm, below=10mm of grup.south, anchor=north] {Async-signal-safe olmayan fonksiyon çağırmak tanımsız davranıştır.};
\end{tikzpicture}
\end{document}
